%%%%%%%%%%%%%%%%%%%%%%%%%%%%%%%%%%%%%%%%%%%%%%%%%%%%%%%%%%%%
%%%%%%%%%%%%%%%%%%%%%%%%%%%%%%%%%%%%%%%%%%%%%%%%%%%%%%%%%%%%
%%%%%%%%%%%%%%%%%%%%%%%%%%%%%%%%%%%%%%%%%%%%%%%%%%%%%%%%%%%%
%%%%%%%%%%%%%%%%%%%%%%%%%%%%%%%%%%%%%%%%%%%%%%%%%%%%%%%%%%%%
\chapter{Introdução ao C}

\setcounter{page}{1}    % set page to 1 again to start arabic count
\pagenumbering{arabic}

O C é uma linguagem {\it case sensitive}, compilada, com tipagem estática e fraca. 
{\it Case sensitive} porque letras maiúsculas são consideradas diferentes de minúsculas. Seu código não roda interativamente como nas linguagens interpretadas. Deve ser compilado, ou seja, gerar código executável antes de rodar. Se diz que o C tem tipagem estática pois o tipo de uma variável, uma vez definido, não muda mais. Se diz que C tem tipagem fraca pois é possível realizar operações entre variáveis de tipos diferentes, como somar inteiros e reais.


%%%%%%%%%%%%%%%%%%%%%%%%%%%%%%%%%%%%%%%%%%%%%%%%%%%%%%%%%%%%
%%%%%%%%%%%%%%%%%%%%%%%%%%%%%%%%%%%%%%%%%%%%%%%%%%%%%%%%%%%%
%%%%%%%%%%%%%%%%%%%%%%%%%%%%%%%%%%%%%%%%%%%%%%%%%%%%%%%%%%%%
\section{Palavras reservadas}

\begin{itemize}

\item Palavras reservadas ou palavras-chave: são palavras que não podem ser usadas para nomear variáveis ou funções criadas pelo usuário.

{\tt
\begin{tabular}{lll}
char     & void     & for     \\
double   & const    & if      \\
float    & switch   & else    \\
int      & case     & typedef \\
long     & default  & struct  \\
short    & break    & return  \\
signed   & continue & main    \\
unsigned & while    & sizeof  \\
static   & do       & enum    \\
\end{tabular}
}

\item A lista de palavras reservadas depende da versão do C. As palavras abaixo, por exemplo, surgiram no padrão C23.

{\tt
\begin{tabular}{llll}
constexpr & bool     & true     & false   \\
\end{tabular}
}



\end{itemize}


%%%%%%%%%%%%%%%%%%%%%%%%%%%%%%
\begin{comment}
\begin{tabular}{ll}
  \toprule
  \textbf{Categoria} & \textbf{Palavra Reservada} \\
  \midrule
    \multirow{4}{*}{Tipos de Dados Primitivos} 
    & \texttt{char} \\
    & \texttt{double} \\
    & \texttt{float} \\
    & \texttt{int} \\
  \midrule
    \multirow{4}{*}{Modificadores de Tipo} 
    & \texttt{long} \\
    & \texttt{short} \\
    & \texttt{signed} \\
    & \texttt{unsigned} \\
  \midrule
    Tipos Especiais & \texttt{void} \\
  \midrule
    Qualificadores de Tipo & \texttt{const} \\
  \midrule
    \multirow{5}{*}{Controle de Fluxo (Seleção)} 
    & \texttt{switch} \\
    & \texttt{case} \\
    & \texttt{default} \\
    & \texttt{if} \\
    & \texttt{else} \\
  \midrule
    \multirow{2}{*}{Controle de Fluxo (Repetição)} 
    & \texttt{for} \\
    & \texttt{while} \\
  \midrule
    \multirow{2}{*}{Controle de Fluxo (Desvio)} 
    & \texttt{break} \\
    & \texttt{continue} \\
  \midrule
    Definição de Tipos & \texttt{typedef} \\
  \midrule
    Tipos Estruturados & \texttt{struct} \\
  \midrule
    Controle de Funções & \texttt{return} \\
  \midrule
    Operadores & \texttt{sizeof} \\
  \midrule
    Ponto de Entrada & \texttt{main} \\
  \bottomrule
\end{tabular}
\end{comment}
%%%%%%%%%%%%%%%%%%%%%%%%%%%%%%


%%%%%%%%%%%%%%%%%%%%%%%%%%%%%%%%%%%%%%%%%%%%%%%%%%%%%%%%%%%%
%%%%%%%%%%%%%%%%%%%%%%%%%%%%%%%%%%%%%%%%%%%%%%%%%%%%%%%%%%%%
%%%%%%%%%%%%%%%%%%%%%%%%%%%%%%%%%%%%%%%%%%%%%%%%%%%%%%%%%%%%
\section{Um programa em C mínimo}


\begin{itemize}

\item Um programa em C mínimo: este programa não faz nada, simplesmente declara uma função \MAIN\ vazia. A função \MAIN\ é o ponto de entrada de qualquer projeto em C, ou seja, por onde o programa começa a execução. Neste exemplo não há nenhuma instrução dentro da \MAIN.

\begin{lstlisting}
int main(void)
{
}
\end{lstlisting}


\item Outro programa em C: este programa simplesmente imprime uma mensagem na tela através de uma chamada à função \PRINTF. Para usar \PRINTF, é necessário incluir o arquivo de cabeçalho {\tt stdio.h}.

\begin{lstlisting}
#include <stdio.h>
int main(void)
{
    printf("Hello world!\n");
}
\end{lstlisting}


\item Ainda outro programa em C: como no exemplo anterior, imprime uma mensagem na tela. Além disso, espera o usuário pressionar Enter antes de terminar a execução através de uma chamada à função \GETCHAR.

\begin{lstlisting}
#include <stdio.h>
int main(void)
{
    printf("Hello world!\n");
    printf("Pressione Enter para continuar...\n");
    getchar();
}
\end{lstlisting}

\end{itemize}


%%%%%%%%%%%%%%%%%%%%%%%%%%%%%%%%%%%%%%%%%%%%%%%%%%%%%%%%%%%%
%%%%%%%%%%%%%%%%%%%%%%%%%%%%%%%%%%%%%%%%%%%%%%%%%%%%%%%%%%%%
%%%%%%%%%%%%%%%%%%%%%%%%%%%%%%%%%%%%%%%%%%%%%%%%%%%%%%%%%%%%
\section{Programando em C}

Para programar em C é necessário um compilador. Se seu sistema operacional for Windows, instale o Orwell Dev-C++, conhecido como {\it devcpp}, disponível em

{\tt http://sourceforge.net/projects/orwelldevcpp/}

Se seu sistema operacional for Linux, instale o gcc rodando o comando

{\tt sudo apt install gcc}

\begin{itemize}

\item Como criar um programa em C

\subitem 1 - Editar o código fonte e salvar em um arquivo com extensão c.
\subitem 2 - Compilar o código fonte usando um compilador.
\subitem 3 - Rodar o programa.

\item Usando o devcpp: o devcpp é um ambiente integrado de programação, com editor e compilador na mesma aplicação.

\subitem 1 - Criar um novo {\it source file} com ctrl + n.
\subitem 2 - Editar e salvar o arquivo com extensão c.
\subitem 3 - Pressionar F9 para compilar e rodar.
\subitem 4 - Ler as mensagens do compilador e corrigir os erros.
\subitem 5 - {\it Go to} 3.

\item Usando o gcc: o gcc é um compilador de linha de comando.

\subitem 1 - Criar um novo {\it source file} com um editor de texto como o emacs.
\subitem 2 - Editar e salvar o arquivo com extensão c.
\subitem 3 - Compilar usando o gcc. Antes de rodar a linha de comando abaixo, substitua o conteúdo {\tt\PLACEHOLDER{destacado}} pelos nomes dos arquivos.

{\tt gcc \PLACEHOLDER{nome\_arq\_fonte.c} -o \PLACEHOLDER{nome\_arq\_executável}}

\subitem 4 - Ler as mensagens do compilador e corrigir os erros de compilação.
\subitem 5 - Chamar o executável caso tenha compilado e corrigir os erros de lógica.
\subitem 6 - {\it Go to} 3.

\item A linguagem C mudou consideravelmente ao longo dos anos, e, por isso, códigos antigos podem resultar em erros e \textit{warnings} quando usamos novas versões do gcc. O Ubuntu 24 usa o gcc 13, que usa {\tt -std=c17} como versão da linguagem. Já o Ubuntu 26 usa o gcc 15, que usa {\tt -std=c23}. Para usar o padrão C23 com o gcc 13, é necessário usar a opção {\tt std=c2x} na linha de comando.

{\tt gcc -std=c2x \PLACEHOLDER{nome\_arq\_fonte.c} -o \PLACEHOLDER{nome\_arq\_executável}}

\subitem Para descobrir em tempo de execução a versão usada na compilação, pode-se imprimir o valor da macro  {\tt \_\_STDC\_VERSION\_\_} como na linha de código a seguir.

\begin{lstlisting}
printf("Standard = %ld\n", __STDC_VERSION__);
\end{lstlisting}

\end{itemize}


%%%%%%%%%%%%%%%%%%%%%%%%%%%%%%%%%%%%%%%%%%%%%%%%%%%%%%%%%%%%
%%%%%%%%%%%%%%%%%%%%%%%%%%%%%%%%%%%%%%%%%%%%%%%%%%%%%%%%%%%%
%%%%%%%%%%%%%%%%%%%%%%%%%%%%%%%%%%%%%%%%%%%%%%%%%%%%%%%%%%%%
\section{Literais}
\label{sec:lit}

Literais são valores com os quais podemos trabalhar. Podem ser de diversos tipos:

\nobreakspace

{\tt
\begin{tabular}{llll}
Números inteiros:        & 0                     & Números reais:    & 2.0       \\
                         & -1                    &                   & 3.1415    \\
                         & 333                   &                   & -1000.0   \\
\\
Inteiros hexadecimais:   & 0xff                  & Inteiros octais:  & 077       \\
                         & 0x002a                &                   & 06        \\
                         & 0xd                   &                   & 02342     \\
\\
Caracteres:              & 'a'                   & Strings:          & "Hello world"       \\
                         & '4'                   &                   & "bom dia"           \\
                         & '\textbackslash n'    &                   & "Idade = \%d anos"   \\
\\
Valores booleanos:       & true                  &                   &                      \\
                         & false                 &                   &                      \\
\end{tabular}
}

\begin{itemize}

\item Para usar os valores booleanos, é necessário incluir o arquivo de cabeçalho {\tt stdbool.h}. No padrão C23 e no C++, esses valores se tornaram palavras-chave e não é necessário incluir o cabeçalho.

\end{itemize}


%%%%%%%%%%%%%%%%%%%%%%%%%%%%%%%%%%%%%%%%%%%%%%%%%%%%%%%%%%%%
%%%%%%%%%%%%%%%%%%%%%%%%%%%%%%%%%%%%%%%%%%%%%%%%%%%%%%%%%%%%
%%%%%%%%%%%%%%%%%%%%%%%%%%%%%%%%%%%%%%%%%%%%%%%%%%%%%%%%%%%%
\section{Chamada de funções}

\begin{itemize}

\item Funções são trechos de código encapsulados que executam alguma operação.
\item Funções são chamadas pelo nome.
\item Recebem zero ou mais parâmetros entre parênteses. O número de parâmetros é definido na sua criação e pode ser constante 
(as funções {\tt sin}, {\tt cos}, {\tt log} recebem um parâmetro; 
a função {\tt pow} recebe dois parâmetros)
ou variável (\PRINTF\ recebe um ou mais parâmetros).
\item Podem retornar um valor de algum tipo de dado. Esse tipo (real, inteiro, caracter etc) é definido no momento da sua criação. O valor retornado pode mudar dependendo dos parâmetros que a função recebe ao ser chamada.
\item Podem também não retornar nada. São as funções de tipo {\tt void}.

\item Exemplos de chamadas:

\begin{lstlisting}
printf("Hello\n");
printf("O mundo termina no ano de %d\n", 2012);
sin(1.0);
pow(2.0, 3.0);
printf("O seno de pi/2 eh %f\n", sin(3.14/2.0));
\end{lstlisting}

\end{itemize}


%%%%%%%%%%%%%%%%%%%%%%%%%%%%%%%%%%%%%%%%%%%%%%%%%%%%%%%%%%%%
%%%%%%%%%%%%%%%%%%%%%%%%%%%%%%%%%%%%%%%%%%%%%%%%%%%%%%%%%%%%
%%%%%%%%%%%%%%%%%%%%%%%%%%%%%%%%%%%%%%%%%%%%%%%%%%%%%%%%%%%%
\section{A biblioteca de funções matemáticas}

Para usá-la, adicione a linha seguinte ao início do código:

\begin{lstlisting}
#include <math.h>
\end{lstlisting}

\begin{itemize}

\item Caso esteja usando Linux, adicionar a opção {\tt -lm} à linha de comando de chamada ao gcc:

{\tt gcc -lm \PLACEHOLDER{nome\_arq\_fonte.c} -o \PLACEHOLDER{nome\_arq\_executável}}


\item Contém funções como

{\tt
\begin{tabular}{lll}
sin      & cos      & tan     \\
asin     & acos     & atan    \\
log      & log10    & pow     \\
ceil     & floor    & fabs    \\
\end{tabular}
}

\item Contém constantes como

{\tt
\begin{tabular}{lll}
M\_PI    & M\_E     &         \\
\end{tabular}
}

\item Abaixo está um exemplo de programa que usa essa biblioteca. O programa imprime uma mensagem com o valor do seno de de pi/4 através de uma chamada a \PRINTF.

\begin{lstlisting}
#include <stdio.h>
#include <math.h>
int main(void)
{
    printf("O seno de pi/4 eh %f\n", sin(3.14/4.0));
}
\end{lstlisting}

\end{itemize}


%%%%%%%%%%%%%%%%%%%%%%%%%%%%%%%%%%%%%%%%%%%%%%%%%%%%%%%%%%%%
%%%%%%%%%%%%%%%%%%%%%%%%%%%%%%%%%%%%%%%%%%%%%%%%%%%%%%%%%%%%
%%%%%%%%%%%%%%%%%%%%%%%%%%%%%%%%%%%%%%%%%%%%%%%%%%%%%%%%%%%%
\section{Comentários}

Comentários são textos ignorados pelo compilador. Dentro de um comentário é possível escrever qualquer coisa, sem risco de erros de compilação. Em C existem dois tipos de comentário:

\begin{itemize}

\item Comentários de uma linha: começa com {\tt //} e vai até o final da linha.
\item Comentários de múltiplas linhas: começa com {\tt /*} e termina com {\tt */}.

\item Exemplo de programa com comentários.

\begin{lstlisting}
/* 
Programa que imprime o valor do seno
de 45 graus, ou pi/4.
*/
#include <stdio.h> // biblioteca do printf
#include <math.h>  // biblioteca do sin
int main(void)
{
    printf("O seno de pi/4 eh %f\n", sin(3.14/4.0));
}
\end{lstlisting}


\end{itemize}


%%%%%%%%%%%%%%%%%%%%%%%%%%%%%%%%%%%%%%%%%%%%%%%%%%%%%%%%%%%%
%%%%%%%%%%%%%%%%%%%%%%%%%%%%%%%%%%%%%%%%%%%%%%%%%%%%%%%%%%%%
%%%%%%%%%%%%%%%%%%%%%%%%%%%%%%%%%%%%%%%%%%%%%%%%%%%%%%%%%%%%
\section{Identificadores}
\label{sec:ident}

Identificadores são os nomes de variáveis, funções, constantes e outros objetos, tanto os criados pelo usuário quanto os contidos em bibliotecas.

\begin{itemize}
\item O primeiro caracter pode ser letra ou sublinhado
\item Os demais caracteres podem ser letras, dígitos ou sublinhado.
\item Não pode ser igual a nenhuma palavra reservada.
\item Exemplos:

{\tt
\begin{tabular}{llll}
Corretos:           & \_ABC                 & Incorretos:     & 123       \\
                    & contador              &                 & i+1       \\
                    & i                     &                 & a!        \\
                    & x2                    &                 & r\$       \\
\end{tabular}
}

\end{itemize}


%%%%%%%%%%%%%%%%%%%%%%%%%%%%%%%%%%%%%%%%%%%%%%%%%%%%%%%%%%%%
%%%%%%%%%%%%%%%%%%%%%%%%%%%%%%%%%%%%%%%%%%%%%%%%%%%%%%%%%%%%
%%%%%%%%%%%%%%%%%%%%%%%%%%%%%%%%%%%%%%%%%%%%%%%%%%%%%%%%%%%%
\section{Variáveis}

Variáveis são posições de memória com um nome, usadas para armazenar valores. Os valores armazenados podem ser modificados a qualquer momento.

\begin{itemize}
\item O nome de uma variável deve ser um identificador válido, como descrito na Seção \ref{sec:ident}.
\item Sua declaração e inicialização devem aparecer no código antes de seu primeiro uso.
\item A declaração de uma variável é onde definimos seu tipo e seu identificador, isto é, seu nome.
\item A inicialização de uma variável é onde atribuímos a ela algum valor.
\item A declaração de uma variável tem o seguinte formato:

{\tt \PLACEHOLDER{tipo de dados} \PLACEHOLDER{identificador};}

\item Exemplos de declarações de variáveis:
\begin{lstlisting}
int x;
int i;
float f;
char c;
unsigned int file_size;
unsigned char idade;
unsigned short int ano;
double p1;
long long int n;
long double value;
\end{lstlisting}


\item Variáveis podem ser
\subitem Locais, quando declaradas dentro de uma função.
\subitem Globais, quando declaradas fora de funções.
\subitem Parâmetros, quando declaradas na lista de parâmetros da definição de uma função.

\item Variáveis locais e globais podem ser inicializadas, ou seja, receber algum valor, no momento de sua declaração. O formato de uma declaração com inicialização é o seguinte:

{\tt \PLACEHOLDER{tipo de dados} \PLACEHOLDER{identificador} = \PLACEHOLDER{valor};}

Lembrando que identificador é o nome. O valor pode ser um literal, como descrito na Seção~\ref{sec:lit}, outra variável, uma chamada de função que retorna algum valor, ou mesmo uma expressão matemática.

\item Exemplos de declarações de variáveis com inicialização:
\begin{lstlisting}
int x = 1;
int i = 0;
float e = 2.7;
char c = 'a';
double p = sin(M_PI/4.0);
double p2 = p;
double p_dup = 2.0 * p;
\end{lstlisting}

\end{itemize}


%%%%%%%%%%%%%%%%%%%%%%%%%%%%%%%%%%%%%%%%%%%%%%%%%%%%%%%%%%%%
%%%%%%%%%%%%%%%%%%%%%%%%%%%%%%%%%%%%%%%%%%%%%%%%%%%%%%%%%%%%
%%%%%%%%%%%%%%%%%%%%%%%%%%%%%%%%%%%%%%%%%%%%%%%%%%%%%%%%%%%%
\section{Constantes}

Como variáveis, constantes são posições de memória que armazenam algum valor. Porém ao contrário de variáveis, o valor não pode ser modificado.

\begin{itemize}

\item A declaração de uma constante tem o seguinte formato:

{\tt const \PLACEHOLDER{tipo de dados} \PLACEHOLDER{identificador} = \PLACEHOLDER{valor};}

\item Exemplos de declarações de constantes:
\begin{lstlisting}
const long int c = 299792458;
const float pi = 3.141592653589793238462;
const float G = 6.67428E-11;
const double h = 6.62606896E-34;
\end{lstlisting}

\item Alternativamente, há a palavra reservada \CONSTEXPR, que garante que o valor é definido em tempo de compilação. Isso é útil, por exemplo, para definir valores usados como um \CASE\ dentro do \SWITCH. Esta palavra está disponível a partir do padrão C23 e também no C++.

{\tt constexpr \PLACEHOLDER{tipo de dados} \PLACEHOLDER{identificador} = \PLACEHOLDER{valor};}

\end{itemize}


%%%%%%%%%%%%%%%%%%%%%%%%%%%%%%%%%%%%%%%%%%%%%%%%%%%%%%%%%%%%
%%%%%%%%%%%%%%%%%%%%%%%%%%%%%%%%%%%%%%%%%%%%%%%%%%%%%%%%%%%%
%%%%%%%%%%%%%%%%%%%%%%%%%%%%%%%%%%%%%%%%%%%%%%%%%%%%%%%%%%%%
\section{Tipos de dados}

Em um computador digital, todo e qualquer dado é armazenado em forma de bits. O conteúdo de uma posição de memória pode ser interpretado como um número inteiro, real ou até mesmo uma string. Portanto, para que um dado seja interpretado corretamente, é necessário associar a esse dado o seu tipo. Na linguagem C os tipos mais importantes são

\begin{itemize}
\item {\tt int}: valor inteiro. Normalmente ocupa 4 bytes, sendo capaz de representar inteiros de -2147483648 a 2147483647.
\item {\tt float}: valor real. Normalmente ocupa 4 bytes, sendo capaz de representar valores reais com cerca de 8 dígitos significativos.
\item {\tt double}: valor real. Normalmente ocupa 8 bytes, sendo capaz de representar valores reais com cerca de 16 dígitos significativos.
\item {\tt char}: valor inteiro pequeno ou caracter. Ocupa 1 byte, sendo capaz de representar inteiros de -128 a 127.
\item {\tt bool}: valor lógico. Pode ser {\tt true} ou {\tt false}. Para usar o tipo \BOOL, é necessário incluir o arquivo de cabeçalho {\tt stdbool.h}. No padrão C23 e no C++, esse tipo se tornou palavra-chave e não é necessário incluir o cabeçalho.

\item {\tt void}: tipo vazio. Usado para declarar funções que não retornam valor ou ponteiros sem tipo associado.
\end{itemize}

Alguns tipos de dados podem ser alterados por modificadores:

\begin{itemize}
\item {\tt unsigned}: pode modificar os tipos {\tt int} e {\tt char}. Uma variável {\tt unsigned} é sempre positiva, e pode representar valores duas vezes maiores que sua correspondente {\tt signed}
\item {\tt signed}: pode modificar os tipos {\tt int} e {\tt char}. 
Por padrão toda variável {\tt int} e {\tt char} é {\tt signed}, ou seja, com sinal.
\item {\tt short}: pode modificar o tipo {\tt int}, diminuindo sua capacidade. Normalmente, em máquinas de 64 bits, um short int ocupa 2 bytes.
\item {\tt long}: pode modificar os tipos {\tt int} e {\tt double}, aumentando sua capacidade. Normalmente, em máquinas de 64 bits, um long int ocupa 8 bytes e um long double, 16 bytes.

\end{itemize}

Para descobrir quanto espaço ocupa uma variável de um certo tipo de dados, usar o operador unário \SIZEOF. \SIZEOF\ pode ser usado com tipos de dados ou com variáveis. No primeiro caso, o uso de parênteses é obrigatório, e no segundo, é opcional. Observe o exemplo abaixo.

\begin{lstlisting}
#include <stdio.h>
int main(void)
{
    char c;
    printf("Um int ocupa %ld bytes\n", sizeof(int));
    printf("A variavel c ocupa %ld bytes\n", sizeof(c));
}
\end{lstlisting}


%%%%%%%%%%%%%%%%%%%%%%%%%%%%%%%%%%%%%%%%%%%%%%%%%%%%%%%%%%%%
%%%%%%%%%%%%%%%%%%%%%%%%%%%%%%%%%%%%%%%%%%%%%%%%%%%%%%%%%%%%
%%%%%%%%%%%%%%%%%%%%%%%%%%%%%%%%%%%%%%%%%%%%%%%%%%%%%%%%%%%%
\section{Expressões numéricas}
\label{sec:expr}

Expressões numéricas são compostas por operadores e operandos. Os operandos podem ser literais, funções que retornam algum valor, variáveis, constantes, ou mesmo outras expressões. Os operadores podem ser de atribuição, de comparação, aritméticos, lógicos ou bit a bit.

\begin{itemize}
\item Uma expressão elementar contém apenas um operando:
\subitem {\tt \PLACEHOLDER{literal}}
\subitem {\tt \PLACEHOLDER{função que retorna valor}}
\subitem {\tt \PLACEHOLDER{variável}}
\subitem {\tt \PLACEHOLDER{constante}}

\item Exemplos de expressões elementares:
\begin{lstlisting}
10
sin(M_PI)
x
M_PI
\end{lstlisting}

\item Uma expressão composta pode conter vários operadores e operandos:
\subitem {\tt ( \PLACEHOLDER{expressão} \PLACEHOLDER{operador binário} \PLACEHOLDER{expressão} )}
\subitem {\tt ( \PLACEHOLDER{operador unário} \PLACEHOLDER{expressão} )}

\item Exemplos de expressões compostas:
\begin{lstlisting}
2 + 2
-1
x = 10
sin(M_PI) + 1.0
y = (3 + x)
n = 2 * (3 + i)
\end{lstlisting}


\end{itemize}

\subsection{O operador de atribuição \emph{=}}

É usado para atribuir um valor a uma variável. É importante salientar que, antes de usar uma variável, ela deve ser inicializada, ou seja, algum valor deve ser atribuído a essa variável. Uma variável não inicializada pode conter um valor imprevisível.

\begin{itemize}
\item A forma geral de uma atribuição é a seguinte:

{\tt \PLACEHOLDER{identificador da variável} = \PLACEHOLDER{expressão};}

Onde é claro que {\tt expressão} pode ser elementar, como um valor literal, outra variável, uma constante, uma chamada a função; ou composta, como uma operação aritmética, lógica etc.

\item Exemplos de atribuições:
\begin{lstlisting}
n = 10;
m = n;
y = - (3 + x);
n = 3 * (4 + i);
\end{lstlisting}

\end{itemize}

\subsection{Operadores aritméticos}

Expressões com operadores aritméticos retornam valores numéricos, inteiros ou reais.

~\\

{\tt
\begin{tabular}{|l|llll|}
\hline
Precedência:        &    &    &     &                                    \\
maior               & ++ & -- &     & (incremento e decremento unário)   \\
                    & -  &    &     & (menos unário)                     \\
                    & *  & /  & \%  & (multiplicaćão, divisão e módulo)  \\
menor               & +  & -  &     & (soma e subtração)                 \\
\hline
\end{tabular}
}

\begin{itemize}

\item Operadores de maior precedência são avaliados antes dos de menor precedência.
\item Operadores com mesmo nível de precedência são avaliados da esquerda para a direita.
\item É possível usar parênteses para modificar a ordem das operações
\item {\tt i++} é o mesmo que {\tt i = i + 1}.
\item {\tt i--} é o mesmo que {\tt i = i - 1}.

\item Nos exemplos abaixo, que valor é atribuído à variável à esquerda do operador {\tt =}?

\begin{lstlisting}
a = 2 * 3 + 1;
b = 1 + 2 * 3;
c = 10 / 2 * 5;
d = 10 % 3;
e = 2 * 10 % 4;
f = 3 * (2 + 1);
g = 2;
h = 3 * ++g;
i = 3 * g++;
\end{lstlisting}

\end{itemize}

\subsection{Operadores lógicos e de comparação}

Expressões com operadores lógicos ou de comparação retornam valores do tipo {\tt bool}, ou seja, {\tt true} ou {\tt false}.

~\\

{\tt
\begin{tabular}{|l|lllll|}
\hline
Precedência:        &      &    &    &    &                                    \\
maior               & !    &    &    &    & (negação lógica)                   \\
                    & >    & >= & <  & <= & (desigualdade)                     \\
                    & ==   & != &    &    & (igualdade e diferença)            \\
                    & \&\& &    &    &    & (and lógico)                       \\
menor               & ||   &    &    &    & (or lógico)                        \\
\hline
\end{tabular}
}

\begin{itemize}

\item Nos exemplos abaixo, que valor é atribuído à variável à esquerda do operador {\tt =}?

\begin{lstlisting}
a = 1 < 2;
b = true || false;
c = true && false;
d = 1 < 2 || 4 < 2;
e = 1 < 2 && 4 < 2;
f = 1 < 1;
g = !(3 < 2);
e = 3 == 4;
e = 3 != 4;
\end{lstlisting}

\item Para testar se o valor de uma variável está em um certo intervalo, a construção é a seguinte: 

\begin{lstlisting}
f > 0 && f < 1
\end{lstlisting}

para testar se f está entre 0 e 1. Jamais omitir o {\tt \&\&} como na expressão INCORRETA abaixo:

\begin{lstlisting}
0 < f < 1
\end{lstlisting}

\end{itemize}

\subsection{Operadores bit a bit}

Operam sobre o valor binário de uma variável. A variável só pode ser do tipo \INT, \CHAR\ ou suas variantes.

~\\

{\tt
\begin{tabular}{|l|llll|}
\hline
Precedência:        &       &    &     &                                    \\
maior               & \LNOT &    &     & (complemento de um)                \\
                    & <<    & >> &     & (shift left e shift right)         \\
                    & \LAND &    &     & (and binário)                      \\
                    & \LXOR &    &     & (xor binário)                      \\
menor               & \LOR  &    &     & (or binário)                       \\
\hline
\end{tabular}
}


\begin{itemize}

\item Nos exemplos abaixo, que valor é atribuído à variável à esquerda do operador {\tt =}?

\begin{lstlisting}
a = 3 | 1;
b = 1 >> 1;
c = 1 << 1;
d = 1 << 2;
e = 5 ^ 3;
f = ~ 0;
\end{lstlisting}

\end{itemize}


%%%%%%%%%%%%%%%%%%%%%%%%%%%%%%%%%%%%%%%%%%%%%%%%%%%%%%%%%%%%
%%%%%%%%%%%%%%%%%%%%%%%%%%%%%%%%%%%%%%%%%%%%%%%%%%%%%%%%%%%%
%%%%%%%%%%%%%%%%%%%%%%%%%%%%%%%%%%%%%%%%%%%%%%%%%%%%%%%%%%%%
\section{Diretivas de pré-processamento}

Diretivas de pré-processamento são usadas pelo compilador para modificar o código-fonte antes de iniciar a compilação propriamente dita. São indicadas pelo símbolo {\tt \#}.

\begin{itemize}

\item Um exemplo já visto anteriormente é o \INCLUDE. A linha abaixo, antes da compilação, é substituída pelo conteúdo do arquivo de cabeçalho {\tt stdio.h}

\begin{lstlisting}
#include <stdio.h>
\end{lstlisting}
 
\item Outro exemplo é o \DEFINE, usado para definir macros. Macros são identificadores que são substituídos pela string que o sucede ao longo do código-fonte. No exemplo abaixo, a primeira linha define {\tt TAM\_MAX} como 10, portanto, antes da compilação, esse identificador será substituído por 10 na declaração de {\tt array}.

{\tt \#define \PLACEHOLDER{identificador} \PLACEHOLDER{string}}

\begin{lstlisting}
#define TAM_MAX 10
int array[TAM_MAX]
\end{lstlisting}

\item O identificador pode ser usado para definir macros com parâmetros, como no exemplo a seguir. Deve-se ter cuidado pois, como a substituição é puramente textual, o valor de {\tt x} pode não ser avaliado corretamente a depender da precedência dos operadores. Isso pode ser evitado com os parênteses ao redor {\tt x}, que garantem que seu valor é calculado antes da multiplicação.

\begin{lstlisting}
#define QUADRADO(x) ((x)*(x))
printf("%d\n", QUADRADO(5+5));
\end{lstlisting}

\end{itemize}


%%%%%%%%%%%%%%%%%%%%%%%%%%%%%%%%%%%%%%%%%%%%%%%%%%%%%%%%%%%%
%%%%%%%%%%%%%%%%%%%%%%%%%%%%%%%%%%%%%%%%%%%%%%%%%%%%%%%%%%%%
%%%%%%%%%%%%%%%%%%%%%%%%%%%%%%%%%%%%%%%%%%%%%%%%%%%%%%%%%%%%
\section{Exercícios}

\begin{enumerate}

  \item Crie uma função \MAIN\ vazia.

  \item Crie uma função \MAIN\ e dentro dela declare as seguintes variáveis:
    \begin{enumerate}
      \item um inteiro chamado {\tt n}.
      \item um real de precisão simples chamado {\tt f}.
      \item um real de precisão dupla chamado {\tt d}.
      \item um caracter chamado {\tt c}.
      \item um inteiro longo chamado {\tt l}.
     \end{enumerate}
          
  \item Em C, as instruções devem ser terminadas com ponto e vírgula. A declaração de uma variável é o primeiro lugar do programa onde a variável aparece e é onde seu tipo de dados é definido. A inicialização é onde algum valor é atribuído a ela pela primeira vez. A declaração e inicialização podem ocorrer na mesma instrução. Por exemplo:
%%%%%%%%%%
\begin{lstlisting}
	int i = 0;
	float x = 5.2;
	float y = x*x;
\end{lstlisting}
%%%%%%%%%%

Crie uma função \MAIN\ e dentro dela declare e inicialize na mesma instrução:
    \begin{enumerate}
      \item um inteiro chamado {\tt val}.
      \item um real de precisao simples chamado {\tt x}.
      \item um real de precisão dupla chamado {\tt xd}.
      \item um caracter chamado {\tt tiny}.
      \item um inteiro longo chamado {\tt huge}.
    \end{enumerate}

  \item Crie uma função \MAIN\ que faz o mesmo que na questão anterior e em seguida, para cada variável, imprime uma mensagem com seu nome, a palavra "vale", e seu valor. Para uma variável chamada pipoca, por exemplo:

	{\tt pipoca vale 3.141592}

  \item Crie uma função \MAIN\ que faz o mesmo que na questão anterior, porém, ao invés de declarar variáveis, declare constantes.

  \item Crie uma função \MAIN\ e dentro dela declare e inicialize na mesma instrução cinco variáveis inteiras. Para inicializá-las use expressões com três operadores aritméticos cada. Na mesma linha, coloque em um comentário o resultado da expressão. Por exemplo:
%%%%%%%%%%
\begin{lstlisting}
	int i = 5 + 2 * 3 * 4; // 29
\end{lstlisting}
%%%%%%%%%%

  \item Crie uma função \MAIN\ e dentro dela declare e inicialize na mesma instrução cinco variáveis inteiras. Para inicializá-las use expressões com cinco operadores aritméticos cada. Na mesma linha, coloque em um comentário o resultado da expressão.

  \item Crie uma função \MAIN\ e dentro dela declare e inicialize na mesma instrução cinco variáveis reais de precisão dupla. Para inicializá-las use expressões com dois operadores aritméticos e uma chamada a função em cada expressão. Na mesma linha, coloque em um comentário o resultado da expressão.

  \item Crie uma função \MAIN\ e dentro dela declare e inicialize na mesma instrução cinco variáveis reais de precisão simples. Para inicializá-las use expressões com três operadores aritméticos cada. Na mesma linha, coloque em um comentário o resultado da expressão.

  \item Crie uma função \MAIN\ que imprime a mensagem \verb|"Hello\n"|.

  \item Crie uma função \MAIN\ que imprime a mensagem \verb|"Hello\n"|, \\\verb|"Pressione Enter para continuar\n"|, e em seguida chama \verb|getchar()|.

  \item O que acontece quando:
    \begin{enumerate}
      \item atribuímos um valor \INT\ a uma variável \FLOAT?
      \item atribuímos um valor \FLOAT\ a uma variável \INT?
    \end{enumerate}

  \item Crie um programa que calcule e imprima a área das figuras planas especificadas. Para um paralelogramo com lados 10 e 4, com ângulo de 60 graus entre eles, por exemplo:
%%%%%%%%%%
\begin{lstlisting}
float a = 10;
float b = 4;
float deg = 60;
float h = b * sin(deg * M_PI / 180);
float area = a * h;
printf("area = %f\n", area);
\end{lstlisting}
%%%%%%%%%%

    \begin{enumerate}
      \item um retângulo com lados 3 e 6.
      \item um quadrado com lado 5.
      \item um triângulo retângulo com catetos 3 e 4.
      \item um triângulo retângulo com cateto igual a 5 e hipotenusa igual a 13.
      \item um triângulo equilátero com lado 4.
      \item um círculo com raio 2.
    \end{enumerate}

  \item Crie um programa que calcule e imprima o volume dos sólidos especificados.
    \begin{enumerate}
      \item um cubo com aresta igual a 5.
      \item um paralelepípedo com lados 3, 4 e 5.
      \item um cilindro cuja base tem raio igual a 3 e altura igual a 4.
      \item uma pirâmide com base quadrada de lado 5 e altura igual a 6.
      \item uma esfera com raio igual a 1.
      \item um prisma cuja base é um triângulo equilátero com lado igual a 2 e altura igual a 10.
    \end{enumerate}

  \item Considere um material com densidade de 0.420 g/cm$^3$. Crie um programa que calcule e imprima a massa em g e kg dos objetos descritos abaixo.
    \begin{enumerate}
      \item uma folha de 19 $\times$ 270 cm e espessura de 0.6 mm.
      \item uma chapa de 210 $\times$ 82 cm e espessura de 3.5 cm.
      \item uma ripa de 1 $\times$ 2 cm e comprimento de 1.2 m.
      \item um sarrafo de 3.5 $\times$ 4 cm e comprimento de 3 m.
      \item um caibro de 5 $\times$ 5 cm e comprimento de 3 m.
      \item uma viga de 4 $\times$ 9 cm e comprimento de 4 m.
      \item um tarugo com diâmetro de 15 mm e comprimento de 60 cm.
    \end{enumerate}

  \item Considere os vetores ${\bf v} = (2, 2, 0)$ e ${\bf w} = (3, -3, 0)$. Crie um programa que calcule e imprima
    \begin{enumerate}
      \item o vetor ${\bf v} + {\bf w}$.
      \item o vetor ${\bf v} - {\bf w}$.
      \item a norma de ${\bf v}$.
      \item a norma de ${\bf w}$.
      \item o coeficiente angular de ${\bf v}$, i.e. o ângulo formado entre o vetor e o eixo $x$.
      \item o coeficiente angular de ${\bf w}$.
      \item o produto escalar de ${\bf v}$ por ${\bf w}$.
      \item o produto vetorial de ${\bf v}$ por ${\bf w}$.
    \end{enumerate}

  \item Considere a equação de segundo grau $x^2 - 3x - 10 = y$. Crie um programa que calcule e imprima
    \begin{enumerate}
      \item o valor de $\Delta = b^2 - 4ac$.
      \item suas duas raízes reais, dadas por $x = (-b \pm \sqrt{\Delta}) / 2a$.
      \item a coordenada $x$ do vértice, dada por $x = -b/2a$.
      \item a coordenada $y$ do vértice, dada por $y = -\Delta/4a$.
    \end{enumerate}

  \item Sabendo que $\log_a b = \log b / \log a$, crie um programa que calcule e imprima
    \begin{enumerate}
      \item o valor de $\log_2 1024$.
      \item o valor de $\log_3 729$.
      \item o valor de $\log_4 256$.
      \item o valor de $\log_5 625$.
    \end{enumerate}

  \item Crie um programa que imprima em uma tabela os valores de 0 a 100 graus Celsius em intervalos de 10 graus na primeira coluna. Na segunda coluna imprima os valores correspondentes em Kelvin e na terceira, os valores correspondentes em Fahrenheit. Seja $c$ o valor em Celsius, $k$ o valor em Kelvin e $f$ o valor em Fahrenheit, considere que $k = c + 273.15$ e $f = c \times 9 / 5 + 32$.

\end{enumerate}
